% ===========================================================================
%  08 -- Conclusion
%
%  Source of truth: Chapters 1, 6, and 7
%  Evidence:       evidence_bank.md entries for RQ1--RQ3
%  Claim boundary: technical auditability and visual-fact reliability are
%                  supported; controlled pedagogical effectiveness is future
%                  work.
% ===========================================================================
\chapter{Conclusion}\label{ch:conclusion}

This thesis began from a practical problem shared by many procedural training
settings: a learner often needs help while acting inside a complex simulated
system, but useful help must be grounded in what is currently true. A fluent
foundation-model response is not enough when the next correct action depends on
telemetry, procedure order, display state, and earlier actions. The central
answer developed in this thesis is that foundation models can be made suitable
for procedural tutoring when they are placed inside an auditable evidence loop.

The F/A-18C cold-start task provided a demanding case for that idea. It required
the tutor to combine telemetry-visible cockpit state, visual evidence from
display pages, retrieved procedural knowledge, manual or out-of-layout
boundaries, and validation before producing overlay or textual guidance. The
result is not a claim that learning effectiveness has already been proven. It is
a technical and methodological contribution: a way to make model-assisted
procedural help state-aware, inspectable, and bounded by available evidence.

% ===========================================================================
\section{Summary of Completed Work}\label{sec:conc-summary}

The completed work answers the three research questions as follows.

\paragraph{RQ1: evidence integration and tutoring architecture.}
The thesis shows how simulator telemetry, retrieved procedural sources, and
visual fact observations can be integrated into a constrained help cycle. The
runtime builds an evidence packet before model generation, infers candidate
procedure steps from declarative gates, asks the model to respond inside a
pack-derived schema, validates or repairs the response through a harness, and
executes only allowlisted overlay actions. Logging and replay then make the
resulting help cycle inspectable after the fact. This answers RQ1 at the level
of technical auditability: the LLM is not the authority for procedure state, but
a generator constrained by evidence, schema, harness, and procedure-pack rules.

\paragraph{RQ2: visual fact extraction and ontology design.}
The thesis also shows that small-data VLM adaptation can support reliable
visual-fact extraction for the cockpit display states required by the tutor. The
13-fact ontology and LoRA training pipeline produce the strongest primary result
on the Qwen3.5-9B line, with approximately 0.99 fact accuracy on two independent
13-fact holdouts and sharply reduced critical false positives relative to the
base model. Supplementary Qwen3.6-27B and Gemma-4-31B results support the same
data-and-ontology recipe within the evaluated backbones. The more general lesson
is representational: the VLM should report locally verifiable visual facts,
while procedural completion should be inferred downstream from visual facts,
telemetry, time, and procedure context.

\paragraph{RQ3: formative VR pilot evidence.}
The completed Quest~3 pilot provides formative evidence about whether the
runtime and instrumentation can be exercised with participants in immersive
simulation. In the reviewed data, the four with-tutor participants completed the
33-step procedure, while the five without-tutor participants did not. Omission
errors disappeared in the tutor condition, while assistance-coupled errors
remained. The help-cycle logs also show that the tutor can operate live with VLM
calls, overlay actions, harness repairs, and replayable exports. These findings
answer RQ3 descriptively. They support feasibility, instrumentation readiness,
and error-pattern analysis; they do not establish causal learning gains,
workload reduction, retention, transfer, or speed improvement.

Together, these answers define the thesis contribution. First, the work provides
an evidence-grounded tutoring runtime pattern for procedural simulation, where
generated help is downstream of explicit state evidence and validation. Second,
it contributes a visual-fact ontology and adaptation method that separates VLM
perception from procedural reasoning. Third, it delivers a pilot-ready
evaluation path for studying grounded tutoring with real VR participants. The
important point across all three is not the presence of an LLM or VLM by itself,
but the boundary placed around foundation-model output.

The scope remains deliberately bounded. The evaluation covers one procedure,
one cockpit layout, local visual holdouts, replay and harness readiness, and a
small non-random pilot. The current evidence supports technical reliability and
auditability more strongly than educational effectiveness. It does not show that
the tutor generalises across domains or that the pilot differences were caused
by the tutor alone. Those boundaries are not weaknesses to hide; they are what
make the contribution inspectable.

% ===========================================================================
\section{Future Work}\label{sec:conc-future}

The immediate next step is a controlled human subject evaluation. Such a study
should keep the current export and replay infrastructure, but add randomised
assignment, balanced participant metadata, complete questionnaire and workload
capture, pre-defined coding rules, and a sample large enough for inferential
analysis. Only that kind of study can test whether, and how, the technically
grounded help cycle changes learning, workload, retention, or transfer.

The technical agenda is equally clear. The procedure-pack and visual-fact
approach should be extended to additional procedures, but each extension needs
new task analysis, evidence-source mapping, visual ontology design, reviewed
data, and independent holdout evaluation. The VLM pipeline should also move
beyond isolated screenshots where procedure state depends on temporal evidence:
short visual histories, repeated confirmation, or multi-frame inputs are natural
next mechanisms for completion-like facts. Finally, the full help cycle should
be evaluated as an integrated system, including latency, overlay correctness,
harness repair frequency, and expert-rated response usefulness.

A longer-term version of the system could learn from in-situ corrections by
instructors or learners. That direction would turn live tutoring into a data
collection loop, but it requires careful safeguards: correction labels may be
noisy, updates must not erase already reliable behaviour, and every adapted
model must still be evaluated against untouched holdouts.

The title of this thesis, \emph{Explain Reality}, captures the final design
commitment. A procedural tutor should not merely explain a checklist, and it
should not merely generate plausible language about a task. It should explain
the learner's current reality: what is observable, which evidence supports it,
which procedure rule applies, what remains uncertain, and why the next response
is allowed. The contribution of this thesis is an architecture and evidence
method for making that explanation auditable.
